\documentclass[11pt,a4paper]{article}
\usepackage[T1]{fontenc}
\usepackage[margin=25mm]{geometry}
\usepackage{amsmath}
\usepackage{listings}
\usepackage{hyperref}
\hypersetup{hidelinks, pdftitle={Aarush A1--A5: Library Seat Management}, pdfauthor={Aarush}}
\lstset{basicstyle=\ttfamily\small, columns=fullflexible,
  keepspaces=true, breaklines=true, frame=single, showstringspaces=false}
\setlength{\parindent}{0pt}
\setlength{\parskip}{6pt}
\newcommand{\routine}[2]{\subsection{\texorpdfstring{\texttt{#1}}{#2}}}
\newcommand{\librarymeaning}{\textbf{Library operation.} }
\newcommand{\arguments}{\textbf{Inputs.} }
\newcommand{\steps}{\textbf{Steps.} }
\newcommand{\result}{\textbf{Result and errors.} }
\newcommand{\example}{\textbf{Example.} }

\title{Smart Library: Seat Management\\Aarush's Tasks A1--A5}
\author{Function-by-function Python implementation guide}
\date{}

\begin{document}
\maketitle

This guide explains \texttt{src/library\_sim/memory.py} through library
operations. A1--A5 here are Aarush's task identifiers in
\texttt{division\_of\_work.txt}; they are not assignment numbers.
The module follows the existing \texttt{library\_sim.pyi} interface contract.

\textbf{Audit snapshot:} checked against code revision \texttt{f69cd13} on
22 November 2026, using Python 3.12.3. The supporting repository report is
\href{../implementation_audit.md}{\texttt{docs/implementation\_audit.md}}.
This is the function guide for the implemented memory module, not the complete
A2 submission README or a claim that the Final project is finished.

\tableofcontents
\newpage

\section{The library model and module boundaries}

Each reader represents an operating-system process. Two librarians represent
two execution units. The six indexed reading seats represent contiguous memory.
A reader requesting several seats is a modelling abstraction: the request
reserves that many adjacent slots, rather than modelling literal furniture use.

\begin{center}
\begin{tabular}{ll}
\textbf{Library operation} & \textbf{Operating-system concept} \\
Find an adjacent group of free seats & Contiguous allocation \\
Reserve seats for a reader & Allocate memory to a process \\
Wait in the seating line & Block for memory admission \\
Wait while seated for a librarian & Wait in the scheduler's ready queue \\
Return seats when finished & Deallocate memory \\
Measure occupied seats over time & Time-weighted memory utilization
\end{tabular}
\end{center}

A seat contains a reader ID, or Python's \texttt{None} if free. Diagrams use
a dot for \texttt{None}. Indices start at zero:
\begin{lstlisting}
Seat index:  0   1   2   3   4   5
Owner:      R1  R1   .   .  R2   .
\end{lstlisting}

The engine owns arrival ordering, lifecycle transitions, admission timestamps,
and librarian scheduling. The memory module only manages reservations,
the admission line, and seating statistics. A seated reader retains the
reservation while waiting for a librarian and, in Final mode, while blocked
for a book. Seats are returned only on completion.

At each timestamp, the A2 engine handles completions and releases, arrivals,
FIFO admission, and librarian dispatch, in that order. It then samples the
final seating state once. All times below are integer virtual time units;
there is no sleeping. Only the coordinator thread accesses this module.

\section{A1: Find and choose seats}

\texttt{FreeRun} means \texttt{tuple[int, int]}: the first seat index and
the group's length. The tuple \texttt{(2, 3)} means seats 2, 3, and 4,
not seats 2 through 3. Groups are maximal: adjacent free seats belong to
one group.

\routine{free\_runs(slots)}{free runs}
\librarymeaning Walk along the seating register and identify every uninterrupted
group of empty seats.

\arguments A sequence of reader IDs and \texttt{None} values. An empty
sequence is allowed.

\steps Remember where a free group begins. When an occupied seat ends the
group, record its start and length. Also record a group that reaches the
end of the room.

\result A new list of maximal free groups, in increasing seat-index order.
The input is not modified. Empty and fully occupied rooms give an empty list.

\example The diagram above gives \texttt{[(2, 2), (5, 1)]}.

\routine{choose\_run(runs, k, allocator)}{choose run}
\librarymeaning Decide where a reader requesting \texttt{k} adjacent seats should
sit. This is a recommendation; it does not reserve seats.

\arguments Valid free groups, a positive request size, and one of the
existing \texttt{Allocator} enum values.

\steps Inspect a copy of the groups in starting-index order. Ignore groups
shorter than the request. First Fit takes the first fitting group; Best Fit
keeps the shortest fitting group; Worst Fit keeps the longest. Equal lengths
retain the lower index. Placement is always at the group's beginning.

\result The chosen starting index, or \texttt{None} if no group fits.
An invalid allocator or nonpositive request raises \texttt{ValueError}.
The supplied groups are unchanged.

\example For one seat in \texttt{X X . . X .}, First Fit returns 2,
Best Fit returns 5, and Worst Fit returns 2. For three seats, all return
\texttt{None}, although three seats in total are free.

\section{A2: SeatPool --- the reservation register}

The room stores a simple seat-owner list and a dictionary mapping each
seated reader to their \texttt{(start, length)} reservation. The dictionary
records exactly what must be returned on completion. No compaction or
movement of existing reservations takes place.

\routine{SeatPool.\_\_init\_\_(allocator, capacity=6)}{SeatPool constructor}
\librarymeaning Open an empty reading room with the chosen seating policy.

\arguments An allocator and a positive seat capacity, normally six.

\steps Store the policy and capacity, fill the seating list with
\texttt{None}, and create an empty reservation dictionary.

\result A new room. Invalid capacity or allocator raises \texttt{ValueError}.
Larger test rooms do not change the standard six-seat requirement.

\example \texttt{SeatPool(Allocator.FIRST\_FIT)} opens six free seats.

\routine{SeatPool.allocate(reader\_id, k)}{SeatPool allocate}
\librarymeaning Reserve an entire adjacent group for a reader.

\arguments A reader ID and a request between 1 and the room's capacity.

\steps Reject an existing reservation for that ID or an invalid size.
Find free groups and ask \texttt{choose\_run} for a starting index. Only
after a fitting group is found, write the ID into all requested seats and
record the reservation.

\result The first reserved seat index. Return \texttt{None} if no group
fits, without changing any seat or recording a partial reservation.
Duplicate reservations and invalid sizes raise \texttt{ValueError}.

\example Allocating two seats to R1 in an empty room returns 0 and makes
seats 0--1 belong to R1. Index 0 is a successful result; callers must check
\texttt{is None}, rather than treating zero as failure.

\routine{SeatPool.release(reader\_id)}{SeatPool release}
\librarymeaning Return exactly the seats reserved by the finishing reader.

\arguments The ID of a seated reader.

\steps Remove their reservation record and set those recorded seats to
\texttt{None}. Leave all other readers' seats untouched. The next free-group
scan naturally joins newly adjacent empty seats.

\result The released \texttt{(start, length)}. A missing reservation,
including a second release, raises \texttt{KeyError} without changing seats.

\example Releasing R2 below joins both neighboring free groups:
\begin{lstlisting}
Before:  .   .  R2  R2   .   .
After:   .   .   .   .   .   .
Returned reservation: (2, 2)
New free groups:      [(0, 6)]
\end{lstlisting}

\routine{SeatPool.free\_runs()}{SeatPool free runs}
\librarymeaning Ask which adjacent groups are currently available in this room.

\arguments No extra arguments; inspect this pool's seats.

\steps Call the pure \texttt{free\_runs} function on the internal register.

\result A separate list in starting-index order; no seating changes.

\example An empty six-seat room returns \texttt{[(0, 6)]}.

\routine{SeatPool.occupied}{SeatPool occupied property}
\librarymeaning Count seats reserved, including seats belonging to readers still
waiting for a librarian.

\arguments A property read, with no extra arguments or parentheses.

\steps Count non-\texttt{None} owners in the seating register.

\result An integer between zero and capacity.

\example \texttt{R1 R1 . . R2 .} has occupancy 3, although only two
distinct readers hold reservations.

\routine{SeatPool.largest\_free\_run}{SeatPool largest free run property}
\librarymeaning Find the largest adjacent group another reader could use.

\arguments A property read of the current room.

\steps Scan current free groups and keep the greatest length.

\result A length, with zero when the room is full. No state changes.

\example \texttt{R1 R1 . . R2 .} gives 2.

\routine{SeatPool.external\_fragmentation()}{SeatPool external fragmentation}
\librarymeaning Measure how much free space lies outside the largest adjacent group.

\arguments No extra arguments; use current occupancy and free groups.

\steps Compute
\[
F_{\mathrm{ext}}=(\text{capacity}-\text{occupied})-
                  \text{largest free group length}.
\]
\result A nonnegative integer; both empty and full rooms give zero.

\example In \texttt{. . R2 R2 . .}, four seats are free but the largest
group has two seats, so fragmentation is $4-2=2$. A three-seat request
cannot fit. This indicator does not mean every scattered seat is unusable:
one- and two-seat requests may still fit.

\routine{SeatPool.snapshot()}{SeatPool snapshot}
\librarymeaning Take a seating photograph for inspection or a later diagram.

\arguments No extra arguments.

\steps Copy the owner list into a tuple.

\result An immutable tuple. Later allocations and releases do not alter
previous photographs.

\example Take a snapshot of an empty room, then seat R1; the saved
snapshot still contains six \texttt{None} values.

\section{A3: AdmissionQueue --- the line for seats}

This queue is distinct from the ready queue. Admission asks whether a reader
can enter the room; the scheduler chooses which admitted reader gets a
librarian. Reader priority does not allow jumping the admission line.

\routine{AdmissionQueue.\_\_init\_\_(pool)}{AdmissionQueue constructor}
\librarymeaning Open a waiting line for a particular reading room.

\arguments The \texttt{SeatPool} that will make reservations.

\steps Keep the supplied pool and create an empty Python list of requests.

\result An empty queue. It shares the real room, rather than copying it.

\example Creating the queue does not reserve any seats.

\routine{AdmissionQueue.enqueue(spec)}{AdmissionQueue enqueue}
\librarymeaning Add an arriving reader to the end of the seating line.

\arguments An existing, validated \texttt{ReaderSpec}; its ID and
\texttt{space\_req} identify the request. The engine supplies arrival order,
with lexicographic ID order for simultaneous arrivals.

\steps Append the reader record. Do not sort by priority or change state.

\result No return value and no seat reservation yet. The input parser is
responsible for valid, unique IDs and feasible request sizes.

\example Enqueuing R1, R2, then R3 preserves that order.

\routine{AdmissionQueue.admit\_ready()}{AdmissionQueue admit ready}
\librarymeaning Seat as many readers from the front as possible without queue jumping.

\arguments No extra arguments; use the attached pool and waiting line.

\steps Ask the pool to allocate for the first reader. After success,
remove that reader and continue. On \texttt{None}, stop immediately without
removing the request or considering later readers.

\result A list of \texttt{(reader ID, first seat)} pairs in admission order.
An empty or blocked queue returns an empty list. Allocation errors propagate,
and the unsuccessful request remains queued. The engine records $M=t$ and
changes admitted readers to READY.

\example R1 requests four seats, R2 three, and R3 one. R1 is seated;
R2 and R3 both wait despite two free seats. After R1 finishes, R2 receives
seats 0--2 and R3 receives seat 3.

\routine{AdmissionQueue.head}{AdmissionQueue head property}
\librarymeaning Read the ID of the reader at the front of the line.

\arguments A property read of the queue.

\steps Inspect the first request if one exists.

\result The ID, or \texttt{None} if empty. The reader is not removed.

\example In the blocked example above, the head is R2.

\routine{AdmissionQueue.waiting\_ids()}{AdmissionQueue waiting IDs}
\librarymeaning Inspect the complete seating line in order.

\arguments No extra arguments.

\steps Build a new list of IDs from the queued reader records.

\result A separate list; changing it cannot change the queue.

\example The blocked line returns \texttt{["R2", "R3"]}.

\routine{AdmissionQueue.\_\_len\_\_()}{AdmissionQueue length}
\librarymeaning Count readers still outside waiting for seats.

\arguments Normally called with \texttt{len(queue)}.

\steps Return the length of the waiting list.

\result An integer count, excluding readers already admitted.

\example With R1 seated and R2/R3 waiting, \texttt{len(queue)} is 2.

\section{A4: Room-usage statistics}

\routine{PoolStats(name, capacity, ...)}{PoolStats constructor}
\librarymeaning Produce an immutable room-usage report. The dataclass supplies
this constructor automatically; there is no separate handwritten initializer.

\arguments All seven fields, in the following order:
\begin{center}
\begin{tabular}{ll}
\texttt{name} & Pool label: \texttt{"SPACE"} \\
\texttt{capacity} & Number of seats \\
\texttt{peak\_occupied} & Largest sampled occupied-seat count \\
\texttt{occupied\_integral} & Total seat-time \\
\texttt{max\_ext\_frag} & Largest sampled fragmentation value \\
\texttt{start\_time} & Beginning of the observation window \\
\texttt{end\_time} & End of that window
\end{tabular}
\end{center}

\steps Store the supplied fields. The tracker supplies consistent values;
the record itself does not repeat input validation.

\result A frozen report. Field assignment raises \texttt{FrozenInstanceError}
from Python's \texttt{dataclasses} module.

\example \texttt{PoolStats("SPACE", 6, 4, 14, 0, 5, 12)} records a peak
of four occupied seats and 14 units of seat-time.

\routine{PoolStats.utilization\_pct}{PoolStats utilization property}
\librarymeaning Calculate what percentage of available seat-time was reserved.

\arguments A property read of a valid report. Let $K$ be capacity and
$E=\text{end time}-\text{start time}$.

\steps For positive elapsed time, calculate
\[
\text{utilization percentage}=
\frac{\text{occupied integral}}{K E}\times100.
\]
\result A floating-point percentage, not a fraction. Zero elapsed time
returns \texttt{0.0} to avoid division by zero.

\example For 14 seat-time units, six seats, and seven time units,
the result is $14/(6\times7)\times100=33.3333\ldots\%$.

\routine{OccupancyTracker.\_\_init\_\_(capacity)}{OccupancyTracker constructor}
\librarymeaning Open a usage log for the reading room.

\arguments A positive room capacity.

\steps Set counts and totals to zero. Leave first and last sample times
unset until an observation arrives.

\result An unsampled tracker. Nonpositive capacity raises \texttt{ValueError}.

\example A tracker created for six seats does not assume the first arrival
is at time zero.

\routine{OccupancyTracker.sample(t, pool)}{OccupancyTracker sample}
\librarymeaning Record the room's state after the engine finishes all events at $t$.

\arguments An integer virtual timestamp and a pool of the tracker's capacity.

\steps Reject a backward time or capacity mismatch before changing the log.
On the first sample, remember the starting time. Otherwise add previous
occupancy multiplied by elapsed time to the integral. Then save today's
occupancy and update the maximum occupancy and fragmentation.

\result No return value. Invalid observations raise \texttt{ValueError}.
Equal timestamps are allowed: they add no elapsed seat-time, and the later
observation becomes the current state. Every supplied observation can update
the peaks; the engine must therefore sample only the final state at each timestamp.

\example If two seats were occupied at $t=5$ and four are occupied at $t=8$,
add $2(8-5)=6$, not $4(8-5)$. New reservations do not apply retroactively.

\routine{OccupancyTracker.finish(end\_time)}{OccupancyTracker finish}
\librarymeaning Produce a room-usage report through a chosen ending time.

\arguments An ending timestamp no earlier than the last sample.

\steps Add the final interval using the most recent occupancy to a local
copy of the integral, then construct \texttt{PoolStats}. Do not move the log's
clock or change its totals.

\result A report. Repeated calls do not double-count. A backward end time
raises \texttt{ValueError}. Without samples, return all-zero usage with both
start and end equal to the supplied ending time.

\example Finishing twice at time 12 returns the same report. Finishing
later can extend the report without changing the underlying observations.

\subsection{Worked occupancy timeline}
\begin{center}
\begin{tabular}{rrrr}
\textbf{Interval} & \textbf{Occupied seats} & \textbf{Duration} & \textbf{Seat-time} \\
$[5,8)$ & 2 & 3 & 6 \\
$[8,10)$ & 4 & 2 & 8 \\
$[10,12)$ & 0 & 2 & 0 \\
\hline
Total & & 7 & 14
\end{tabular}
\end{center}
The window starts at 5, not 0. Peak occupancy is four and utilization is
$14/42\times100\%$. In a successful simulation the window runs from first
arrival to last completion. Unsampled intermediate states do not affect peaks.

\section{A5: Tests and integration evidence}

The tests in \texttt{tests/contract/test\_memory\_contract.py} use small,
hand-calculated examples. No external test framework is required.

\routine{reader(reader\_id, seats)}{Test reader helper}
\librarymeaning Prepare a reader's seating request for a queue test.

\arguments An ID and a seat count.

\steps Construct the existing \texttt{ReaderSpec} with arrival zero,
priority one, service one, category REGULAR, and the requested seats.

\result A reader specification. This helper neither seats the reader nor
runs a simulation; it keeps test examples readable.

\subsection{What each test group establishes}
\begin{description}
\item[AllocationFunctions] Empty/full/boundary free groups, input preservation,
the specification's FF/BF/WF choices, distinct policy choices in a larger test
room, tie-breaking, exact fit, no fit, and invalid arguments.
\item[SeatReservations] Six-seat defaults, full-capacity and one-seat
requests, placement by policy, fragmented failures with no partial reservation,
exact release, coalescing on both sides, duplicate/invalid reservations,
double release, configurable capacity, and immutable snapshots.
\item[AdmissionLine] Empty queue behavior, caller-supplied ordering,
independent inspection lists, blocked heads, retained requests, and admission
of every fitting reader after a release.
\item[OccupancyStatistics] Hand-calculated integrals, initial/ending idle
intervals, first arrival after zero, last-interval accounting, repeatable reports,
equal timestamps, sampled peaks, fragmentation peaks, zero duration, frozen
reports, full utilization, and rejected observations without log corruption.
\item[MemoryWithExistingEngine] A hand-calculated FIFO example under both
schedulers and all three allocators, final seat conservation, reader termination,
and equality with the resource-free Final engine baseline.
\end{description}

\subsection{Worked engine check}
All three readers arrive at zero, with equal priority:
\begin{center}
\begin{tabular}{lrrrrr}
\textbf{Reader} & \textbf{Seats} & \textbf{Service} & \textbf{Admission} & \textbf{Start} & \textbf{Completion} \\
R1 & 4 & 4 & 0 & 0 & 4 \\
R2 & 3 & 2 & 4 & 4 & 6 \\
R3 & 1 & 1 & 4 & 4 & 5
\end{tabular}
\end{center}
R3 waits behind R2 until R1 releases its seats. Occupancy is four during
$[0,4)$, four during $[4,5)$, and three during $[5,6)$.
Thus seat-time is $4\times4+4\times1+3\times1=23$, peak occupancy is four,
fragmentation is zero, and utilization is $23/36\times100\%$.

Separate A1 regression tests use the Python golden-output files. The
engine-versus-A1 cross-check itself compares timestamps and summary metrics
against the Python A1 implementation on 11 bundled workloads plus 200 seeded
random workloads under
both schedulers. A room large enough for the sum of all readers' requests
cannot delay admission, so this is a sufficient unlimited-memory baseline
without allocating a million test seats. Golden output contents and scheduling
assertions are unchanged.

At this revision, all 25 memory checks pass. Root discovery finds 134 tests:
133 pass and one missing-module check is skipped because the Final module
is present. The standalone Python A1 suite passes 17 tests and its golden-output
runner passes 20 cases. These are Python 3.12.3 results, not guarantees for
every input or interpreter.

\subsection{Running the checks}
Run these commands from the repository root:
\begin{lstlisting}[language=bash]
PYTHONDONTWRITEBYTECODE=1 python3 -m unittest \
  tests.contract.test_memory_contract -v
PYTHONDONTWRITEBYTECODE=1 python3 -m unittest \
  tests.contract.test_pyi_conformance -v
PYTHONDONTWRITEBYTECODE=1 python3 -m unittest \
  discover -s tests -t . -v
\end{lstlisting}
For the standalone Python A1 baseline, change to
\texttt{Assignment\_1-OS\_Lab\_Python} and run:
\begin{lstlisting}[language=bash]
PYTHONDONTWRITEBYTECODE=1 python3 -m unittest discover -s tests -v
\end{lstlisting}

\section{Design choices and remaining work}

Explicit loops and basic Python containers make the algorithms easy to trace
in the viva. The implementation favors clarity over search-speed optimization.
It implements the existing planned public API without changing its signatures
or the interface contract.

Strict FIFO admission improves arrival-order fairness but can leave usable
seats idle behind a blocked head. Fragmentation can prevent a large request
even when enough seats are free in total. Neither allocation compaction nor
bypassing the head is part of this implementation.

\subsection{Current task status}
This delivery covers Aarush A1--A5. A6 still needs separate A2 inputs and
hand-verified expected outputs. A7/A7b still need the experiment harness and
measured comparisons. Memory contract tests exist, but the separate A7c fake
does not. A8's runtime memory map/timeline is not implemented; the diagrams
here are explanatory examples. A9/A10 still need the complete A2 README,
comparison results, and submission instructions.

\subsection{Integration and packaging now present}
The shared engine and CLI work with real memory. Samar's resource, worker,
trace, and deadlock modules now exist and pass the bundled integrated cases.
Controlled normal-worker overlap, several trace/parser details, and complete
policy comparisons remain unfinished; consult the audit for limitations.

The ZIP builder produces executable A2 and Final archives. A2 excludes
Final-only modules and runs independently. Both archives currently contain
a generated entrypoint, Python package, and one-line placeholder README.
They omit full documentation, samples, and tests. The packaged default
deadlock demo lacks its input. Packaging exists, while complete submission
assembly remains pending.

The deprecated C++ implementation is not required to run the code or tests.

\end{document}
